% ECONOMICS_PROBLEM: {"id": "E31-576", "title": "Phillips-curve nonlinearities", "jel_code": "E31", "source_id": "OP-0242"}
% !TeX program = xelatex
\documentclass[11pt,a4paper]{article}
\usepackage[margin=23mm]{geometry}
\usepackage{fontspec}
\setmainfont{Latin Modern Roman}
\usepackage{amsmath,amssymb,mathtools}
\usepackage{xcolor,enumitem,booktabs,longtable,array,xurl,hyperref}
\definecolor{muted}{HTML}{53636C}
\hypersetup{colorlinks=true,urlcolor=blue,linkcolor=blue}
\setlength{\parindent}{0pt}
\setlength{\parskip}{5pt}
\newcommand{\R}{\mathbb R}
\newcommand{\E}{\mathbb E}
\newcommand{\N}{\mathbb N}
\newcommand{\ind}{\mathbf 1}
\newcommand{\Var}{\operatorname{Var}}
\newcommand{\Cov}{\operatorname{Cov}}
\newcommand{\supp}{\operatorname{supp}}
\DeclareMathOperator*{\argmin}{arg\,min}
\DeclareMathOperator*{\argmax}{arg\,max}
\newcommand{\entrymeta}[3]{{\sffamily\small\color{muted}\textbf{\texttt{#1}}\quad|\quad #2\par #3\par}}
\newcommand{\Problemref}[1]{\texttt{#1}}
\newcommand{\JELref}[2]{\texttt{#2} (JEL \texttt{#1})}
\begin{document}
\section*{Phillips-curve nonlinearities}
\textbf{Problem ID:} \texttt{E31-576}\quad \textbf{JEL:} \texttt{E31}\par
% BEGIN ECONOMICS STATEMENT
\label{problem:OP-0242}
{\sffamily\small\color{muted}Original subject group: Inflation and monetary policy\par}
\label{definitions:problem:30007523}
\textbf{Audit classification:} Published research agenda. Phillips-curve determinants are listed; the source does not specifically assert the kink model is unresolved. Verification concerns the cited version, not first priority or proof that this exact editorial specification remains unsolved on October 8, 2026.
\subsection*{Definitions and precise research target}
\paragraph{Editorial mathematical specification.} Let \(\pi_t\) be quarterly inflation, \(E_t^s\pi_{t+1}\) a specified survey measure of expectations, and \(x_t\) an explicitly constructed measure of real economic slack. Model the conditional inflation relation as
\[
\pi_t=\alpha+\beta E_t^s\pi_{t+1}
+\kappa_1 x_t+\kappa_2(x_t-\tau)_+
+\gamma'w_t+u_t,
\]
where \((r)_+=\max\{r,0\}\), \(w_t\) contains predetermined supply controls, \(\tau\) is a possible threshold, and \(u_t\) contains unobserved cost disturbances. Specify instruments \(Z_t\) satisfying \(E[Z_tu_t]=0\), relevance conditions, and the sample population before estimation.
The target is the identified set for \((\kappa_1,\kappa_2,\tau)\), allowing measurement error in slack and expectations and weak threshold identification when \(\kappa_2\) is near zero. Determine whether a slope change survives changes in the slack measure, supply controls, and monetary-policy regime. Compare a threshold model with smooth nonlinear alternatives on held-out observations; do not infer a kink solely from extreme inflation episodes that also experienced supply disruptions. State the invariance assumptions under which an estimated curve predicts inflation after a policy intervention. The proposed empirical agenda concerns the existence, location, and policy relevance of nonlinearities in a defined setting, rather than whether inflation and activity ever covary.

Take \(x_t,w_t,Z_t\) as observed population random vectors with finite moments, bound \(\tau\) to an interval \([\tau_L,\tau_U]\), and fix a compact coefficient set. Define \(u_t(\vartheta)\) as inflation minus the displayed right-hand-side prediction, and \(\mathcal I(P)=\{\vartheta:E_P[Z_tu_t(\vartheta)]=0\}\). Instruments relevant for the linear slack coefficient need not identify the nonlinear term or threshold. If \(\kappa_2=0\), every \(\tau\) gives the same conditional equation and its location is unidentified. With latent true slack \(x_t^*\) and observed \(x_t=x_t^*+v_t\), add a measurement model or validation data; substituting observed slack does not preserve instrument exogeneity automatically. A descriptive conditional curve does not become a causal policy equation without structural invariance.
\subsection*{Variants and assumption changes}
The following are \emph{editorial variants}, not additional published open conjectures.
\begin{enumerate}
\item \emph{Smooth curve.} Replace the kink by \(f(x)\) in a fixed smoothness ball and normalize \(f(0)=0\). Identify \(f'\) and test whether its sign and curvature are stable across policy regimes.
\item \emph{State-dependent slope.} Let \(\kappa\) depend on a predetermined inflation regime rather than contemporaneous slack. Distinguish a regime-dependent slope from a slack threshold with jointly valid instruments and compare their confidence sets.
\end{enumerate}
\subsection*{References and source audit}
\begin{itemize}
\item European Central Bank. \emph{Forecasting and business cycle analysis}. Undated. Locator: Understanding the inflation process, first selected area (Phillips-curve determinants); no specific kink claim. Primary text checked. \url{https://www.ecb.europa.eu/pub/economic-research/research_agenda/forecasting/html/index.en.html}.
\end{itemize}
Phillips-curve determinants are listed; the source does not specifically assert the kink model is unresolved.
\begin{enumerate}\raggedright
\item\hypertarget{definitions:source:30007523:1}{} European Central Bank. Undated / date unverified. \emph{Forecasting and business cycle analysis}. Understanding the inflation process, first selected area (Phillips-curve determinants); no specific kink claim.
\url{https://www.ecb.europa.eu/pub/economic-research/research_agenda/forecasting/html/index.en.html}.
\end{enumerate}

\subsection*{Common conventions: Source mathematical conventions}

These conventions apply to each entry unless explicitly replaced.

\textbf{Models and identification.} A primitive model \(m\) specifies all probability laws, feasible choices, information, equilibrium and measurement rules used by its target. The admissible class \(\mathcal M\) is fixed independently of outcomes. If \(P_O(m)\) is its observable probability law and \(\tau(m)\) the target, its identified set at observable law \(P\) is
\[
 \mathcal I_\tau(P)=\{\tau(m):m\in\mathcal M,\ P_O(m)=P\}.
\]
Point identification means that this set is a singleton. An empty set signals incompatibility of data and assumptions. Coordinate bounds need not describe the joint set. Bounds are sharp only if every retained target is attainable or arbitrarily approachable by compatible models. Finite bounds require explicit restrictions on unobserved quantities; unrestricted confounding, hidden wealth, extrapolation or arbitrary equilibrium selection can yield uninformative sets.

\textbf{Causal interventions.} A potential outcome \(Y_i(p)\) is the outcome under a complete policy regime \(p\), including timing, financing, information and market spillovers. The target population and units are fixed before treatment. With interference, use the complete assignment vector or a justified exposure mapping; individual no-interference notation alone cannot identify an economy-wide intervention. A difference of mean potential outcomes does not require the cross-world joint distribution. Conditioning on post-treatment migration, survival, enrollment or employment changes the estimand and needs separate justification.

\textbf{Statistical statements.} An estimator, confidence set, forecast or test specifies a sampling law, model class and loss/coverage criterion. Uniform \((1-\alpha)\)-coverage means \(\inf_{m\in\mathcal M}\Pr_m\{\tau(m)\in C_n\}\geq1-\alpha\), or an explicitly asymptotic version. The whole parameter space trivially has coverage, so informativeness requires a stated width, risk or power objective. A forecasting improvement is relative to a named benchmark, information set and evaluation population.

\textbf{Equilibrium and optimization.} An equilibrium comprises feasible plans, optimality, beliefs where needed, budget constraints and all market/resource clearing. The primitives or a stated benchmark define each correspondence \(\mathcal E(p;m)\). Nonemptiness is assumed or proved, never inferred just from compact policy sets. A selected-equilibrium target uses a declared selection rule; a robust target quantifies over all permitted equilibria. Use suprema or infima unless attainment is established. An \(\varepsilon\)-optimal policy is feasible and within \(\varepsilon\) of the finite optimum value. Report an empty feasible set separately.

\textbf{Welfare and accounting.} Normative utility, interpersonal normalization, social weights, numeraire, discounting and the use of public receipts are primitives. Behavioral utility need not coincide with the welfare criterion. Household welfare includes ownership income and rents once; adding profits again to compensating variation generally double-counts them. A monetary equivalent is defined by an explicit transfer and price convention, with monotonicity and range conditions. Average effects, mechanism contributions, transfers and welfare losses are different targets. Infinite sums require convergence and dynamic feasible plans require terminal or no-Ponzi conditions.

\textbf{Source versions.} A working-paper date, revision date and journal publication year are distinguished. A DOI for a journal version is not automatically the DOI of an earlier manuscript. Locators refer to the stated version and printed pages where specified. A metadata-only check does not verify a claim about a full-text conclusion. Every entry's source subsection narrows the provenance claim accordingly.
% END ECONOMICS STATEMENT
\end{document}
